%!TEX encoding = UTF8
%!TEX root = 0-notes.tex

\chapter{Forme échelonnée}
\label{chap:forme-echelonnee}

\section{Forme échelonnée}

\subsection{Opérations élémentaires}

% Type I 
% 
% def permutation de [n]
% def matrice P_sigma e_i = e_sigma-1(i) 
% donc P_sigma (lignes) = sigma(lignes)
% inverse P_sigma = P_sigma-1 = P_sigma^T
% 
% Type II
% matrices diagonales D_i(lambda) = (1, .., 1, lambda, 1, .., 1)
% donc D_i(lambda) (lignes) = lignes avec i-ème x lambda
% D_i(lambda)^-1 = D_i(1/lambda)
% 
% Type III
% G_ij(lambda) = I + lambda en (j, i)
% donc G_ij(lambda)(lignes) = lignes avec j-ième = j-ième + lambda x i-ème
% G_ij(lambda)^-1 = G_ij(-lambda)


\subsection{Méthode du pivot et forme échelonnée réduite}

% Soit A (m x n). On veut créer B inversible tq BA soit le plus simple possible
% exemples matrices simples
% def échelonnée
% def échelonne réduite
% thm : existe B inversible tq BA soit échelonnée réduite

\section{Espaces et sous-espaces vectoriels}

% corollaire : A est inversible si et seulement si elle est carrée (m=n) et sa forme échelonnée réduite est l'identité
% corollaire-prop du chapitre 2 : si k vecteurs de Rn sont L.I., alors k<=n.
% thm : A est inversible ssi carrée et ses colonnes sont L.I. ssi carrée et ses lignes sont L.I.
% 
% famille génératrice
% expression Im(A) = { Ax }_x
% expression Vect(v_i) = span(v_i).
% base : génératrice + L.I.
% base de Rn a exactement n vecteurs

\prop{}{
	Soit $M\in\R^{n \times n}$ une matrice où $n\in\N^*$.
	Supposons que $M$ admette un inverse à gauche : $M^{-G} M = I$.
	
	Alors l'inverse à gauche est unique et est également un inverse à droite : $M^{-G} = M^{-D} = M^{-1}$.
}{prop:inv2}

\pf{}{
	todo
}

\section{Résolution de système linéaire : matrice augmentée}

% Ax = b lorsque A est inversible donne une unique solutoin <=> colonnes de A forment une base
% mais en général peut-être que A a des colonnes redondantes ou pas assez de colonnes, et x est peut-être dans le span de A
%
% echelonnage de (A | b) donne
% B(A | b) = (C | Bb) = (C|b') échelonnée réduite
% Donc A x = b <=> Cx = b'
% thm : si C a une ligne nulle alors que b' n'est pas nul, aucune solution. Sinon, solutions selon la dimension de C (nb variables) - (nb contraintes).
% remarque : on s'en fiche totalement de B en fait.
% exemple Kressner ou autre
% https://www.epfl.ch/labs/anchp/wp-content/uploads/2018/05/ch3.pdf




\section*{Exercices supplémentaires}
\addcontentsline{toc}{section}{Exercices supplémentaires}

\shipoutExercise

%\section*{Solutions}
%\addcontentsline{toc}{section}{Solutions}
%
%\shipoutAnswer



